\documentclass[../../main.tex]{subfiles}
\begin{document}

{\bf [解答]}

% TODO: y=x/9, y=sin \pi x/6のグラフ

与えられた関数を
\begin{align}
  f(x)=\dfrac{x}{9}-\sin\dfrac{\pi x}{6} \label{eq:1}
\end{align}
とおくと，一階微分は
\begin{align}
  f'(x)
  & = \dfrac{1}{9} - \dfrac{\pi}{6}\cos\dfrac{\pi x}{6} \\
  & = \dfrac{6}{\pi}\left[\dfrac{\pi}{54}-\cos\dfrac{\pi x}{6}\right] \label{eq:2}
\end{align}
である．

まず$6\le x$の時$f(x)=0$の解が存在しないことを示す．

\subparagraph{(i)$6\le x\le12$の時}

この時は\cref{eq:1}の第1,2項について
\begin{align}
  \sin\dfrac{\pi x}{6}\le0<\dfrac{1}{9}x
\end{align}
だから，$f(x)>0$である．よって解はない．

\subparagraph{(ii)$12\le x$の時}

この時は$\sin$の最大値が$1$であることより
\begin{align}
  \dfrac{x}{9} \ge \dfrac{12}{9} > 1 \ge \sin\dfrac{\pi x}{6}
\end{align}
だから$f(x)>0$である．よって解はない．
\casesend

以上から$6 \le x $の時は$f(x)=0$の解はない．
そこで$x < 6$で考える．

% TODO: y=f(x)の[3,6]での増減表

この時は\cref{eq:2}より$f(x)$の増減表は上表のようになる．
よって$f(x)$は$3\le x\le6$で単調増加であり，
\begin{align}
  f(4.5)=\dfrac{1-\sqrt{2}}{2}<0, \quad
  f(5)=\dfrac{1}{18}>0
\end{align}
である．よって中間値の定理から$4.5<x<5$に最大の解を持つ．
よって求める整数は$5$である．$\cdots$(答)
\newpage
\end{document}
